\section{Introduction}

Quantum Bogosort is commonly associated with an extreme observer-selection intuition: unfavorable branches are assigned lower first-person accessibility, so conditional experience becomes concentrated on favorable branches. The technical problem is that this intuition mixes at least three distinct objects: the base branch measure, the policy-dependent trajectory map, and the observer-indexed conditioning rule.

This paper separates those objects. Recognition of a QBS-type rule is treated as a causal variable that can change policy. A changed policy can alter ordinary downstream outcomes even before any observer-selection term is introduced. Separately, an accessibility function can reweight first-person conditioning over the policy-dependent histories that remain in the model. This yields a decomposition between ordinary causal policy effects and observer-indexed selection effects.

A second question is where favorable outcome--accessibility alignment comes from. We formalize an adaptive-agent route in which an internal signal predicts the conditional mean of future outcomes. In the clean score-measurable case, accessibility that respects the same ordering inherits a nonnegative outcome covariance. The framework is then generalized beyond deterministic accessibility: total covariance separates alignment of conditional expected outcome and conditional expected accessibility from residual outcome--accessibility dependence that remains after conditioning on the score. This residual term is kept explicit rather than assumed away.

The contribution is therefore not the claim that an external random process becomes objectively lucky, nor the elementary fact that weighted conditioning can change a conditional mean. The contribution is the structured recognition-dependent framework linking policy changes, downstream trajectories, predictive alignment, observer-indexed accessibility, residual dependence, and statistically auditable sufficient conditions.

The central physical uncertainty is explicitly isolated: the mathematical framework does not derive an Everettian accessibility rule from unitary dynamics, decoherence, observer theory, or the Born rule. Any Everett interpretation requires an additional bridge assumption. This separation allows the mathematical, statistical, and decision-theoretic structure to be evaluated independently of the physical interpretation.

\paragraph{Contributions.} The paper provides: (i) exact covariance and tail identities for first-person weighting; (ii) a monotone-accessibility sufficient condition for first-order stochastic dominance; (iii) an exact recognition decomposition into trajectory and conditioning effects; (iv) a policy--QBS interaction theorem and adaptive-rescue corollary; (v) a predictive-calibration mechanism for endogenous outcome--accessibility alignment and a general-accessibility decomposition retaining residual conditional covariance; (vi) classical experiments showing endogenous predictive ordering, paired recognition effects, interaction structure, and branch-wide policy coherence; and (vii) explicit falsification conditions, finite-sample validation layers, and physical-interpretation boundaries.
